\scanpage{50}
\mspage{14}
\margin{\deleted{\unsure{formulation fausse}\illegible{1}}}
``\uline{Der 6'}'' \inserted{(provisoire)} Si $\mathcal X$ est fibré sur
\[\underline\Psi^{\mathrm{op}}=\underline\Phi=\left(\diagram{\begin{tikzcd}&a\arrow[dl]\\1&\\&b\arrow[ul]\end{tikzcd}}\right)\]
(donc donné par, à $\underline Q$-équiv. près, par un diagramme $\diagram{\begin{tikzcd}&X_0\\\unsure{X_a|X_1}\arrow[ur,no head]\arrow[dr]&\\&X_b\end{tikzcd}}$ dans Cat) alors pour $\xi\in\underline D(\mathcal X)$ constant (voire l.c.)
\[\deleted{\unsure{H\xi}}\quad\psi_!(\sigma)\in\underline D(\underline Q)\]
est non seulement cocartésien (sans mérites), mais aussi cartésien.
\margin{\uline{Axiome des fibrations}}
\[\underline\Phi\subset\underline Q^{\mathrm{op}}=\diagram{\begin{tikzcd}[row sep=large,column sep=small,execute at end picture={\draw (\tikzcdmatrixname-2-1.center) ++(.45,0) ellipse (.95cm and 1.35cm);}]&a\arrow[dl]&\\1&&0\arrow[ul]\arrow[dl]\\&b\arrow[ul]&\end{tikzcd}}\]
\uline{NB} Der 6' pour $D$ \deleted{\unsure{implique}} n'est autre que Der 6 pour $D^{\mathrm o}$. Je ne vois pas comment l'un impliquerait l'autre, en dehors du cas exceptionnel où il y aurait \deleted{\unsure{comme}\ (\unsure{\forall\xi})} un foncteur dualisant dans $\mathcal A=D(e)$, et une corresp. donnée par dualité entre cohomologie et homologie.

Finalement, il semblerait que le bon axiome, qui englobe Der 6 et Der 6 bis, soit le suivant :
\margin{\uuline{La définition} \deleted{\unsure{abstraite}}}
\margin{``\uuline{l'axiome géométrique}''}
\uline{Der 6} \deleted{\unsure{bis}\illegible{1}} Si $\sigma\in\underline D(\Delta^1\times\Delta^1)$, $\sigma$ est cartésien ssi il est cocartésien.
\inserted{\uline{De plus} (\uline{à moins} que ça \unsure{ne soit absurde}) :}
Le diagramme squelettique associé est
\mspage{49}

